\subsection{Estado del Arte}

\subsubsection{Arquitecturas híbridas BaaS y FaaS en aplicaciones web modernas}

La práctica contemporánea de desarrollo de aplicaciones web ha consolidado un modelo en el que múltiples servicios gestionados se combinan para soportar distintas responsabilidades del sistema, en lugar de descansar en un único proveedor o tecnología \cite{Newman2021}. En este modelo, los servicios de \textit{Backend as a Service} (BaaS) suelen utilizarse para responsabilidades horizontales bien delimitadas —autenticación, persistencia, almacenamiento—, mientras que las funciones \textit{serverless} bajo el modelo \textit{Function as a Service} (FaaS) se reservan para lógica de negocio específica, integraciones con APIs externas y procesamiento asíncrono. Esta combinación, conocida como \textit{arquitectura híbrida}, ofrece un balance pragmático entre velocidad de desarrollo (heredada del modelo BaaS) y control arquitectónico (aportado por la flexibilidad del modelo FaaS).

\textcite{Michael2021} describe esta tensión en el contexto del crecimiento de \textit{startups}: Firebase ofrece una arquitectura de dos capas en la que las aplicaciones cliente interactúan directamente con servicios \textit{backend}, lo que facilita el desarrollo rápido de productos mínimos viables pero introduce desafíos de escalabilidad, control y seguridad a medida que el producto madura. Frente a este escenario, una vía habitual en la industria consiste en complementar Firebase con servicios de AWS —en particular Lambda, API Gateway, SQS, SNS y CloudWatch—, manteniendo Firebase para las responsabilidades donde sigue siendo idóneo (autenticación, base de datos documental con sincronización en tiempo real) e introduciendo AWS donde se requieren capacidades adicionales (lógica del lado del servidor controlada, observabilidad estructurada, procesamiento asíncrono escalable). Esta arquitectura híbrida es precisamente la adoptada por Paralegales.

\subsubsection{Patrón Backend-for-Frontend: origen, evolución y adopción}

El patrón \textit{Backend-for-Frontend} (BFF) tiene su origen documentado en la experiencia de SoundCloud durante la transición de un monolito hacia microservicios: el equipo encontró que exponer una API genérica para múltiples clientes (web, iOS, Android) generaba acoplamiento entre ellos y degradaba la experiencia de cada uno. La solución fue construir un \textit{backend} específico por cada \textit{frontend}, optimizado para sus necesidades \cite{NewmanBFF}. \textcite{Newman2021} formaliza posteriormente el patrón en su análisis sistemático de microservicios, describiéndolo como una variante del patrón ``API Gateway'' donde, en lugar de un único gateway genérico, existe un gateway dedicado por cada tipo de cliente.

La adopción del BFF se ha extendido a empresas como Netflix, donde se utiliza para adaptar la entrega de contenido a la diversidad de dispositivos cliente, y a frameworks contemporáneos como Nuxt~3 y Next.js, que ofrecen un entorno de ejecución del lado del servidor estrechamente integrado con la capa de presentación. En estos frameworks, el BFF se materializa de manera natural mediante rutas de servidor que actúan como proxy, agregador o coordinador de llamadas a servicios subyacentes. En Paralegales, el BFF implementado sobre Nuxt~3 cumple varias funciones simultáneas: reubica lógica sensible fuera del cliente, encapsula la comunicación con servicios externos no documentados (Rama Judicial, SAMAI), centraliza el manejo de errores y autenticación, y constituye un punto natural de instrumentación para la telemetría del sistema.

\subsubsection{Evolución arquitectónica incremental: Strangler Fig Pattern}

Cuando un sistema en producción presenta problemas estructurales pero no puede detenerse para una reescritura completa, la industria ha consolidado un patrón de evolución incremental conocido como \textit{Strangler Fig Pattern}, descrito por \textcite{Fowler2004Strangler}. La metáfora proviene de las higueras estranguladoras de los bosques tropicales, que crecen alrededor de un árbol huésped hasta reemplazarlo completamente sin que el ecosistema circundante perciba la transición. Aplicado al software, el patrón consiste en introducir gradualmente nuevos componentes que progresivamente asumen responsabilidades del sistema existente, hasta que este último puede ser retirado sin disrupciones.

Este enfoque resulta particularmente apropiado para sistemas con base de usuarios activa, donde una migración total acarrearía riesgos operativos significativos: pérdida de datos, interrupciones del servicio, regresiones funcionales no detectadas y costos de oportunidad asociados al freezing de \textit{features} durante la transición. El \textit{Strangler Fig Pattern} permite distribuir esos riesgos a lo largo del tiempo y validar cada paso de la transformación frente a usuarios reales. La intervención descrita en este trabajo se inscribe en esta filosofía: en lugar de migrar la totalidad de Paralegales desde Firebase hacia AWS en un único movimiento, se introducen capacidades complementarias (BFF, microservicios serverless, observabilidad, IaC) que rodean el sistema existente y absorben progresivamente responsabilidades, preservando la opción de evolución futura sin imponer una sustitución inmediata del proveedor incumbente. El patrón \textit{Handler--Service--Repository} adoptado en el \textit{backend} es, de hecho, una aplicación directa de esta filosofía: aísla el motor de persistencia detrás de una abstracción que permitiría, en el futuro, sustituir Firestore por otro proveedor sin tocar la lógica de negocio.

\subsubsection{Observabilidad en arquitecturas serverless distribuidas}

La adopción de arquitecturas distribuidas basadas en microservicios y funciones \textit{serverless} introduce desafíos particulares para el diagnóstico operativo: una sola acción del usuario puede atravesar el BFF, varios microservicios independientes y múltiples proveedores externos antes de retornar una respuesta. Los enfoques tradicionales de monitoreo —basados en revisar manualmente \textit{logs} servidor por servidor— resultan insuficientes ante esta complejidad. La práctica contemporánea articula tres pilares complementarios: \textit{logs} estructurados, métricas agregadas y trazas distribuidas \cite{Sridharan2018}.

Los \textit{logs} estructurados (típicamente en formato JSON) sustituyen a los registros de texto plano y permiten su indexación, búsqueda y agregación automática. Las métricas derivadas a partir de patrones en los \textit{logs} habilitan paneles de control en tiempo real y alarmas proactivas sin requerir instrumentación adicional. Las trazas distribuidas, finalmente, correlacionan eventos a través de toda la malla de servicios mediante identificadores propagados de petición. En el ecosistema AWS, estos tres pilares se materializan mediante CloudWatch Logs (con \textit{Logs Insights} para consultas relacionales), CloudWatch Metrics (con filtros de métricas deducidas a partir de \textit{logs}) y AWS X-Ray (para trazas distribuidas). La práctica de \textit{Site Reliability Engineering} \cite{Beyer2016} consolida estos elementos en un programa operativo que incluye definición de SLOs, gestión de \textit{error budgets} y reducción sistemática del MTTR. En Paralegales, el subsistema de observabilidad implementado sigue directamente esta línea: interceptores HTTP que emiten eventos JSON estandarizados, métricas deducidas, alarmas vinculadas a notificaciones por Telegram y dashboards desplegados mediante CloudFormation.

\subsubsection{Prácticas SDLC modernas y aseguramiento de calidad continuo}

La capacidad de un equipo para sostener la velocidad de cambio sin degradar la confiabilidad del sistema depende en buena medida de la madurez de sus prácticas de ciclo de vida de desarrollo. \textcite{Humble2010} introdujeron el concepto de \textit{Continuous Delivery} (CD), que plantea automatizar de extremo a extremo el camino entre el \textit{commit} del desarrollador y el despliegue en producción, de modo que cada cambio pueda liberarse de forma rápida, repetible y de bajo riesgo. La CD descansa sobre tres pilares operativos: control de versiones disciplinado, automatización de pruebas a múltiples niveles, y \textit{pipelines} de integración continua que validan cada cambio frente a un conjunto explícito de criterios.

La automatización a nivel local —mediante \textit{hooks} de Git que ejecutan \textit{linters}, formateadores y validadores de mensajes de \textit{commit}— complementa los \textit{pipelines} centralizados al detectar problemas antes de que el código entre al repositorio, reduciendo la fricción de las revisiones y el costo de las regresiones. La estrategia de pruebas, por su parte, se articula clásicamente como una pirámide en cuya base se sitúan numerosas pruebas unitarias rápidas, en el medio pruebas de integración menos numerosas, y en la cúspide un número reducido de pruebas \textit{end-to-end} que validan flujos completos desde la perspectiva del usuario. En Paralegales, esta estrategia se materializa mediante Lefthook como gestor de \textit{hooks}, ESLint y Prettier para análisis estático y formateo, Commitlint con \textit{Conventional Commits} para la trazabilidad del historial, GitHub Actions para los \textit{pipelines} de CI, y Vitest acompañado de la suite de emuladores de Firebase para pruebas unitarias de las reglas de seguridad de Firestore.

\subsubsection{Migración total a AWS: alternativa evaluada y descartada}

Durante la fase exploratoria del proyecto se consideró la alternativa de una migración total de la plataforma desde Firebase hacia AWS, sustituyendo cada servicio de Firebase por su contraparte en AWS: Firebase Authentication por Amazon Cognito, Cloud Firestore por Amazon DynamoDB, Firebase Cloud Functions por AWS Lambda y Firebase Cloud Messaging por Amazon SNS \cite{Michael2021}. La literatura ofrece guías técnicas detalladas para acometer este tipo de transición: \textcite{Hall2017} describe estrategias para migrar usuarios desde Firebase Authentication hacia Cognito (importación por lotes con restablecimiento obligatorio de contraseñas o migración progresiva al iniciar sesión), \textcite{Shank2021} provee un repositorio de código abierto basado en AWS CDK que automatiza la replicación de los servicios clave de Firebase en AWS, y \textcite{Kansara2024} presenta un marco general para procesos ETL hacia bases de datos en la nube apoyado en AWS Step Functions, Lambda y SQS. \textcite{Kyadasu2025} complementan esta perspectiva con prácticas DevOps específicas para migraciones entre proveedores de nube.

A partir de esta evaluación se inició formalmente el trabajo de migración del motor de persistencia desde Firestore hacia DynamoDB, lo que motivó originalmente la introducción del patrón \textit{Handler--Service--Repository} en el \textit{backend}. Sin embargo, tras analizar el balance entre el costo total de la migración —técnico, operativo y de oportunidad— y los beneficios esperados frente a una plataforma ya en producción con base de usuarios activa, la organización decidió descartar la migración total y consolidar en su lugar una estrategia de evolución arquitectónica híbrida. El patrón de repositorios introducido durante esa fase exploratoria se conservó por los beneficios independientes de testabilidad y desacoplamiento que aporta, y deja la puerta abierta a una eventual sustitución futura del motor de persistencia sin acoplar dicha decisión al resto del sistema.
